
\documentclass{article}
\usepackage{colortbl}
\usepackage{makecell}
\usepackage{multirow}
\usepackage{supertabular}

\begin{document}

\newcounter{utterance}

\centering \large Interaction Transcript for game `wordle\_withcritic', experiment `medium\_frequency\_words\_clue\_with\_critic', episode 3 with refineonly{-}Qwen3.5{-}9B.
\vspace{24pt}

{ \footnotesize  \setcounter{utterance}{1}
\setlength{\tabcolsep}{0pt}
\begin{supertabular}{c@{$\;$}|p{.15\linewidth}@{}p{.15\linewidth}p{.15\linewidth}p{.15\linewidth}p{.15\linewidth}p{.15\linewidth}}
    \# & $\;$A & \multicolumn{4}{c}{Game Master} & $\;\:$B\\
    \hline

    \theutterance \stepcounter{utterance}  
    & & \multicolumn{4}{p{0.6\linewidth}}{
        \cellcolor[rgb]{0.9,0.9,0.9}{
            \makecell[{{p{\linewidth}}}]{
                \texttt{\tiny{[P1$\langle$GM]}}
                \texttt{You are a language wizard who likes to guess words by using the given rules.} \\
\\ 
\texttt{Welcome to Wordle! You have six attempts to guess the target word, a valid English word of five lowercase letters (a{-}z). Please use the tags "explanation:" and "guess:" to provide a concise explanation for each guess.} \\
\\ 
\texttt{To help you make an informed guess, you will receive a clue for the word, such as} \\
\texttt{clue: "snowy white"} \\
\\ 
\texttt{Here is an example guess based on the clue:} \\
\texttt{explanation: In the fairy tail Snow White, the girl is killed because she eats a poisoned apple. And the word apple has 5 letters.} \\
\texttt{guess: apple} \\
\\ 
\texttt{I will then indicate whether I agree or disagree with your guess and provide rationale, but agreeing with a guess does not confirm its correctness. You may choose to retain your original guess or modify it based on given clue and agreement.} \\
\\ 
\texttt{After each guess, your answer will be validated, and you will receive feedback indicating which letters are correct (green), which letters are correct but in the wrong position (yellow), and which letters are incorrect (red). This feedback can be useful in determining which letters to include or exclude in your next guess.} \\
\\ 
\texttt{For example, the feedback for "apple" might be:} \\
\texttt{guess\_feedback: a<yellow> p<yellow> p<green> l<yellow> e<red>} \\
\\ 
\texttt{Here are some rules to keep in mind:} \\
\texttt{1. Please do not generate any guess feedback or clue on your own} \\
\texttt{2. Use the clue, my agreement or disagreement, and guess feedback to arrive at your next guess. This should be reflected in your explanation as well.} \\
\\ 
\texttt{Do not generate any other text.} \\
\\ 
\texttt{Let's begin with your first guess.} \\
\\ 
\texttt{clue: sweetheart} \\
            }
        }
    }
    & \\ \\

    \theutterance \stepcounter{utterance}  
    & \multicolumn{4}{p{0.6\linewidth}}{
        \cellcolor[rgb]{0.95,0.95,0.95}{
            \makecell[{{p{\linewidth}}}]{
                \texttt{\tiny{[P2$\rightarrow$P1]}}
                \texttt{{[}DRAFT{]}} \\
\texttt{explanation: The clue "sweetheart" suggests a word related to affection or a partner. "SWEET" is a direct synonym for the first part of the clue and is a common 5{-}letter word.} \\
\texttt{guess: sweet} \\
            }
        }
    }
    & & \\ \\

    \theutterance \stepcounter{utterance}  
    & \multicolumn{4}{p{0.6\linewidth}}{
        \cellcolor[rgb]{0.9,0.9,0.9}{
            \makecell[{{p{\linewidth}}}]{
                \texttt{\tiny{[P1$\rangle$GM]}}
                \texttt{explanation: The clue "sweetheart" suggests a word related to affection or a partner. "SWEET" is a direct synonym for the first part of the clue and is a common 5{-}letter word.} \\
\texttt{guess: sweet} \\
            }
        }
    }
    & & \\ \\

    \theutterance \stepcounter{utterance}  
    & & \multicolumn{4}{p{0.6\linewidth}}{
        \cellcolor[rgb]{0.9,0.9,0.9}{
            \makecell[{{p{\linewidth}}}]{
                \texttt{\tiny{[GM$\rangle$P2]}}
                \texttt{I need your assistance with a word game in which we need to find a 5{-}letter word using a clue, a guess and an explanation for the guess.} \\
\\ 
\texttt{For example, if the clue is "snow white" and the guess is "apple", you will get the information like this:} \\
\texttt{clue: snow white} \\
\texttt{explanation: the fairy tale Snow White involves the girl eating a poisoned apple, which has 5 letters.} \\
\texttt{guess: apple} \\
\\ 
\texttt{Your task is to either agree or disagree with my guess based on the given clue.} \\
\\ 
\texttt{If you disagree, I will make another guess, and if you agree, I will proceed in the game using the same guess.} \\
\\ 
\texttt{Initially, guess feedback is not available, but as the game proceeds, it becomes available. The guess feedback is not for current guess, but rather an indication of what letters can be present in the current guess. A letter highlighted in green indicates that it is a correct letter in the correct position, while yellow indicates that it is a correct letter in the wrong position, and red indicates a wrong letter. At this point, you should use all the provided information, including the clue, guess, explanation, and guess feedback, to decide whether the given guess matches the clue, and also stick to guess feedback or not.} \\
\\ 
\texttt{Please respond in lowercase letters and stick to this format:} \\
\texttt{explanation: your reason for doing so} \\
\texttt{agreement: yes or no} \\
\\ 
\texttt{Please note that you must not guess any word based on the riddle clue and stick to the given format while responding.} \\
\\ 
\texttt{Do not generate any other text.} \\
\\ 
\texttt{Let's begin!} \\
\\ 
\texttt{clue: sweetheart} \\
\texttt{explanation: The clue "sweetheart" suggests a word related to affection or a partner. "SWEET" is a direct synonym for the first part of the clue and is a common 5{-}letter word.} \\
\texttt{guess: sweet} \\
            }
        }
    }
    & \\ \\

    \theutterance \stepcounter{utterance}  
    & & & \multicolumn{4}{p{0.6\linewidth}}{
        \cellcolor[rgb]{0.95,0.95,0.95}{
            \makecell[{{p{\linewidth}}}]{
                \texttt{\tiny{[P1$\rightarrow$P2]}}
                \texttt{{[}DRAFT{]}} \\
\texttt{explanation: the word "sweet" is a direct synonym for the first part of the clue "sweetheart" and fits the 5{-}letter requirement.} \\
\texttt{agreement: yes} \\
            }
        }
    }
     \\ \\

    \theutterance \stepcounter{utterance}  
    & & & \multicolumn{4}{p{0.6\linewidth}}{
        \cellcolor[rgb]{0.9,0.9,0.9}{
            \makecell[{{p{\linewidth}}}]{
                \texttt{\tiny{[GM$\langle$P2]}}
                \texttt{explanation: the word "sweet" is a direct synonym for the first part of the clue "sweetheart" and fits the 5{-}letter requirement.} \\
\texttt{agreement: yes} \\
            }
        }
    }
     \\ \\

    \theutterance \stepcounter{utterance}  
    & & \multicolumn{4}{p{0.6\linewidth}}{
        \cellcolor[rgb]{0.9,0.9,0.9}{
            \makecell[{{p{\linewidth}}}]{
                \texttt{\tiny{[P1$\langle$GM]}}
                \texttt{clue: sweetheart} \\
\texttt{guess\_agreement: yes} \\
\texttt{agreement\_explanation: the word "sweet" is a direct synonym for the first part of the clue "sweetheart" and fits the 5{-}letter requirement.} \\
\\ 
\texttt{Provide your response only in this format.} \\
\texttt{explanation: details} \\
\texttt{guess: word} \\
            }
        }
    }
    & \\ \\

    \theutterance \stepcounter{utterance}  
    & \multicolumn{4}{p{0.6\linewidth}}{
        \cellcolor[rgb]{0.95,0.95,0.95}{
            \makecell[{{p{\linewidth}}}]{
                \texttt{\tiny{[P2$\rightarrow$P1]}}
                \texttt{{[}DRAFT{]}} \\
\texttt{explanation: The user agreed that "sweet" is a valid guess based on the clue "sweetheart". Since the agreement does not confirm the word is correct, I will retain this guess to see the feedback on the letters S, W, E, E, T.} \\
\texttt{guess: sweet} \\
            }
        }
    }
    & & \\ \\

    \theutterance \stepcounter{utterance}  
    & \multicolumn{4}{p{0.6\linewidth}}{
        \cellcolor[rgb]{0.9,0.9,0.9}{
            \makecell[{{p{\linewidth}}}]{
                \texttt{\tiny{[P1$\rangle$GM]}}
                \texttt{explanation: The user agreed that "sweet" is a valid guess based on the clue "sweetheart", but agreement does not confirm correctness. Retaining the guess is necessary to receive feedback on the positions of S, W, E, and T.} \\
\texttt{guess: sweet} \\
            }
        }
    }
    & & \\ \\

    \theutterance \stepcounter{utterance}  
    & & & \multicolumn{2}{p{0.3\linewidth}}{
        \cellcolor[rgb]{0.9,0.9,0.9}{
            \makecell[{{p{\linewidth}}}]{
                \texttt{\tiny{[GM$|$GM]}}
                \texttt{attempts = 1} \\
\texttt{target\_word = flame} \\
\texttt{target\_word\_clue = sweetheart} \\
\texttt{guess = sweet} \\
\texttt{guess\_feedback = s<red> w<red> e<yellow> e<red> t<red>} \\
            }
        }
    }
    & & \\ \\

    \theutterance \stepcounter{utterance}  
    & & \multicolumn{4}{p{0.6\linewidth}}{
        \cellcolor[rgb]{0.9,0.9,0.9}{
            \makecell[{{p{\linewidth}}}]{
                \texttt{\tiny{[P1$\langle$GM]}}
                \texttt{guess\_feedback: s<red> w<red> e<yellow> e<red> t<red>} \\
\\ 
\texttt{Provide your response only in this format.} \\
\texttt{explanation: details} \\
\texttt{guess: word} \\
            }
        }
    }
    & \\ \\

    \theutterance \stepcounter{utterance}  
    & \multicolumn{4}{p{0.6\linewidth}}{
        \cellcolor[rgb]{0.95,0.95,0.95}{
            \makecell[{{p{\linewidth}}}]{
                \texttt{\tiny{[P2$\rightarrow$P1]}}
                \texttt{{[}DRAFT{]}} \\
\texttt{explanation: The feedback indicates that 'E' is present but in the wrong position (yellow), while S, W, and T are incorrect (red). Since the clue is "sweetheart", the target word likely contains 'E' and relates to love or affection. Words like "LOVER" or "BELOVED" come to mind, but "LOVER" has no 'E'. "BELOVED" is 7 letters. "SWEET" had 'E' in the middle. Let's try "LOVES" (plural/verb) or "CARES"? No, clue is sweetheart. Maybe "LOVER" is wrong. How about "SWEET" was a trap? Wait, 'E' is yellow. Maybe the word is "LOVER" but I need an 'E'. "LOVER" has no E. "BELOVED" is too long. "SWEETHEART" is too long. What 5 letter word means sweetheart and has an E? "LOVER" no. "SWEET" no. Maybe "CARE"? No. "LOVE" is 4. "LOVES" has E at end. If E is yellow in sweet (pos 3,4), maybe it's not at end. "LOVER" is the most direct synonym. But no E. Is there another word? "SWEET" {-}> S,W,E,E,T. E is yellow. So E is in the word. Maybe "LOVER" is not it. "SWEET" {-}> maybe "SWEET" is not the word. What about "SWEET" {-}> "SWEET" is 5 letters. Maybe the word is "SWEET" but the user said E is yellow. So S,W,T are out. E is in. What 5 letter word means sweetheart and has E? "LOVER" no. "BELOVED" no. "SWEET" no. Maybe "CARE" no. "LOVE" no. "SWEET" {-}> maybe "SWEET" is not the word. Wait, "SWEET" {-}> S,W,E,E,T. E is yellow. So E is in the word. Maybe the word is "LOVER" but I need an E. "LOVER" has no E. "BELOVED" is 7. "SWEETHEART" is 9. What about "SWEET" {-}> "SWEET" is 5. Maybe the word is "SWEET" but the user said E is yellow. So S,W,T are out. E is in. What 5 letter word means sweetheart and has E? "LOVER" no. "BELOVED" no. "SWEET" no. Maybe "CARE" no. "LOVE" no. "SWE} \\
            }
        }
    }
    & & \\ \\

    \theutterance \stepcounter{utterance}  
    & \multicolumn{4}{p{0.6\linewidth}}{
        \cellcolor[rgb]{0.9,0.9,0.9}{
            \makecell[{{p{\linewidth}}}]{
                \texttt{\tiny{[P1$\rangle$GM]}}
                \texttt{explanation: The feedback indicates 'E' is present but in the wrong position (yellow), while S, W, and T are incorrect (red). Since the clue is "sweetheart", the target word likely contains 'E' and relates to love or affection. Words like "LOVER" or "BELOVED" come to mind, but "LOVER" has no 'E'. "BELOVED" is 7 letters. "SWEET" had 'E' in the middle. Let's try "LOVES" (plural/verb) or "CARES"? No, clue is sweetheart. Maybe "LOVER" is wrong. How about "SWEET" was a trap? Wait, 'E' is yellow. Maybe the word is "LOVER" but I need an 'E'. "LOVER" has no E. "BELOVED" is too long. "SWEETHEART" is too long. What 5 letter word means sweetheart and has an E? "LOVER" no. "SWEET" no. Maybe "CARE"? No. "LOVE" is 4. "LOVES" has E at end. If E is yellow in sweet (pos 3,4), maybe it's not at end. "LOVER" is the most direct synonym. But no E. Is there another word? "SWEET" {-}> S,W,E,E,T. E is yellow. So E is in the word. Maybe "LOVER" is not it. "BELOVED" is too long. "SWEET" no. Maybe "CARE" no. "LOVE" no. "SWEET" {-}> maybe "SWEET" is not the word. Wait, "SWEET" {-}> S,W,E,E,T. E is yellow. So E is in the word. Maybe the word is "LOVER" but I need an 'E'. "LOVER" has no E. "BELOVED" is 7. "SWEETHEART" is 9. What about "SWEET" {-}> "SWEET" is 5. Maybe the word is "SWEET" but the user said E is yellow. So S,W,T are out. E is in. What 5 letter word means sweetheart and has E? "LOVER" no. "BELOVED" no. "SWEET" no. Maybe "CARE" no. "LOVE" no. "SWEET" {-}> maybe "SWEET" is not the word. Wait, "SWEET" {-}> S,W,E,E,T. E is yellow. So E is in the word. Maybe the word is "LOVER" but I need an 'E'. "LOVER" has no E. "BELOVED" is 7. "SWEETHEART"} \\
            }
        }
    }
    & & \\ \\

    \theutterance \stepcounter{utterance}  
    & & & \multicolumn{2}{p{0.3\linewidth}}{
        \cellcolor[rgb]{0.9,0.9,0.9}{
            \makecell[{{p{\linewidth}}}]{
                \texttt{\tiny{[GM$|$GM]}}
                \texttt{Error: The response should contain the 'guess:' keyword exactly once.} \\
            }
        }
    }
    & & \\ \\

    \theutterance \stepcounter{utterance}  
    & & & \multicolumn{2}{p{0.3\linewidth}}{
        \cellcolor[rgb]{0.9,0.9,0.9}{
            \makecell[{{p{\linewidth}}}]{
                \texttt{\tiny{[GM$|$GM]}}
                \texttt{game\_result = ABORT} \\
            }
        }
    }
    & & \\ \\

\end{supertabular}
}

\end{document}
